\section{Fibonacci stabilization on words}\label{sec:stabilization}

This section fixes the local arithmetic object of the paper. Stabilization is defined by reducing Fibonacci valuations modulo $F_{m+2}$ and returning the unique admissible representative. The resulting map $\Fold_m$ is genuinely many-to-one on isolated words, and its residue fibers are the arithmetic input for the threshold and branch-locus theorems proved later.

\subsection{Admissible language and Fibonacci valuation}\label{sec:stab:golden}

Let $(F_n)_{n\ge 1}$ be the Fibonacci sequence with $F_1=F_2=1$ and $F_{n+2}=F_{n+1}+F_n$.

\begin{definition}[Golden-mean legal language]\label{def:Xm}
For $m\ge 1$, let
\[
X_m := \Bigl\{x=(x_1,\dots,x_m)\in\{0,1\}^m : x_i x_{i+1}=0\ \text{for all }1\le i<m\Bigr\}.
\]
\end{definition}

\begin{definition}[Fibonacci valuation]
For $\omega=(\omega_1,\dots,\omega_m)\in\{0,1\}^m$, define
\[
\Val_m(\omega):=\sum_{k=1}^m \omega_k F_{k+1}.
\]
\end{definition}

\begin{proposition}[Finite-window Zeckendorf range]\label{prop:finite-window-range}
For every $m\ge 1$, the restriction
\[
\Val_m|_{X_m}:X_m\longrightarrow \{0,1,\dots,F_{m+2}-1\}
\]
is a bijection.
\end{proposition}

\begin{proof}
By Zeckendorf's theorem, every integer $N\ge 0$ admits a unique expansion
\[
N=\sum_{k\ge 1}\varepsilon_k F_{k+1},\qquad \varepsilon_k\in\{0,1\},\qquad \varepsilon_k\varepsilon_{k+1}=0
\]
\cite{Zeckendorf1972,Lothaire2002}. If $0\le N < F_{m+2}$, then the largest Fibonacci weight that can appear is at most $F_{m+1}$, so the Zeckendorf expansion of $N$ has length at most $m$ and defines a unique element of $X_m$. Thus $\Val_m|_{X_m}$ is surjective.

To prove injectivity, let $x,y\in X_m$ with $\Val_m(x)=\Val_m(y)$. Then $x$ and $y$ are two Zeckendorf expansions of the same integer using weights up to $F_{m+1}$, hence $x=y$ by uniqueness.

Equivalently, admissible words of length $m$ realize each residue class $0,1,\dots,F_{m+2}-1$ exactly once.
\end{proof}

\begin{remark}
Proposition~\ref{prop:finite-window-range} is the finite-window form of Zeckendorf uniqueness that makes a canonical stabilization map possible without any ambiguity at the right boundary. In that sense it is the rigid finite-window analogue of the finite-automaton normalization viewpoint developed for Fibonacci numeration in \cite{Frougny1992,Berstel2001}.
\end{remark}

\subsection{Canonical finite-window stabilization}\label{sec:stab:fold}

\begin{definition}[Canonical stabilization map]\label{def:fold}
For $m\ge 1$, define
\[
\Fold_m:\{0,1\}^m\longrightarrow X_m
\]
by the rule
\[
\Fold_m(\omega)
:=
\bigl(\Val_m|_{X_m}\bigr)^{-1}\!\Bigl(\Val_m(\omega)\pmod{F_{m+2}}\Bigr),
\]
where the residue is taken in $\{0,1,\dots,F_{m+2}-1\}$.
\end{definition}

\begin{remark}
The map $\Fold_m$ is the unique finite-window canonicalization compatible with Fibonacci numeration: compute the raw Fibonacci-weighted value $\Val_m(\omega)$, reduce it modulo $F_{m+2}$, and write the residue in its unique length-$m$ Zeckendorf form. Local carry rewrites such as $11\mapsto 100$ are one implementation of this residue-class stabilization.
\end{remark}

\begin{theorem}[Basic properties of $\Fold_m$]\label{thm:fold-suite}
For every $m\ge 1$, the stabilization map $\Fold_m:\{0,1\}^m\to X_m$ satisfies:
\begin{enumerate}[label=(\roman*),leftmargin=*,itemsep=2pt]
  \item \textbf{Well-definedness.} $\Fold_m$ is well-defined.
  \item \textbf{Idempotence on admissible words.} If $x\in X_m$, then $\Fold_m(x)=x$.
  \item \textbf{Surjectivity.} $\Fold_m$ is onto $X_m$.
  \item \textbf{Fiber partition.} The fibers
  \[
  \mathcal{F}_m(x):=\Fold_m^{-1}(\{x\}),\qquad x\in X_m,
  \]
  form a partition of $\{0,1\}^m$. In particular,
  \[
  \sum_{x\in X_m} d_m(x)=2^m,
  \qquad
  d_m(x):=|\mathcal{F}_m(x)|.
  \]
\end{enumerate}
\end{theorem}

\begin{proof}
(i) follows from Proposition~\ref{prop:finite-window-range}, which gives a unique admissible representative for each residue class modulo $F_{m+2}$.

(ii) If $x\in X_m$, then $\Val_m(x)\in\{0,\dots,F_{m+2}-1\}$ and $\bigl(\Val_m|_{X_m}\bigr)^{-1}(\Val_m(x))=x$.

(iii) Every $x\in X_m$ is a fixed point of $\Fold_m$, so $x$ lies in the image.

(iv) Since $\Fold_m$ is a function on the finite set $\{0,1\}^m$, its inverse images form a partition. Summing the cardinalities of the fibers gives the cardinality of the domain.
\end{proof}

\subsection{Residue fibers and Fourier formula}\label{sec:stab:fiber}

\begin{theorem}[Exact residue characterization of fibers]\label{thm:fiber-residue}
Let $x\in X_m$ and write
\[
r_x:=\Val_m(x)\in\{0,\dots,F_{m+2}-1\}.
\]
Then
\[
\mathcal{F}_m(x)
=
\Bigl\{\omega\in\{0,1\}^m:
\Val_m(\omega)\equiv r_x \pmod{F_{m+2}}
\Bigr\}.
\]
Equivalently,
\[
d_m(x)
=
\#\Bigl\{\omega\in\{0,1\}^m:
\sum_{k=1}^m \omega_k F_{k+1}\equiv r_x \pmod{F_{m+2}}
\Bigr\}.
\]
\end{theorem}

\begin{proof}
By definition,
\[
\Fold_m(\omega)=x
\iff
\bigl(\Val_m|_{X_m}\bigr)^{-1}\!\Bigl(\Val_m(\omega)\pmod{F_{m+2}}\Bigr)=x.
\]
Since $\Val_m|_{X_m}$ is bijective, this is equivalent to
\[
\Val_m(\omega)\equiv \Val_m(x)\pmod{F_{m+2}}.
\]
This proves the claim.
\end{proof}

\begin{remark}[Worked example at $m=5$]
At $m=5$ the admissible language $X_5$ has $F_7=13$ words. The raw word $\omega=11010$ has valuation
\[
\Val_5(11010)=F_2+F_3+F_5=1+2+5=8.
\]
The unique admissible word of length five with residue $8$ modulo $13$ is $00001$, since $\Val_5(00001)=F_6=8$. Thus
\[
\Fold_5(11010)=00001.
\]
In fact
\[
\mathcal F_5(00001)=\{00001,\ 00110,\ 11010\}.
\]
This is the basic local ambiguity studied throughout the paper: isolated windows can fold together even though Section~\ref{sec:stab:conjugacy} will show that overlap removes the ambiguity once $m\ge 3$.
\end{remark}

\begin{theorem}[Discrete Fourier formula for the fiber cardinalities]\label{thm:fiber-dft}
Let $\zeta:=e^{2\pi i/F_{m+2}}$. For every $x\in X_m$ with $r_x=\Val_m(x)$,
\begin{equation}\label{eq:fiber-dft}
d_m(x)
=
\frac{1}{F_{m+2}}
\sum_{j=0}^{F_{m+2}-1}
\zeta^{-j r_x}
\prod_{k=1}^m \bigl(1+\zeta^{jF_{k+1}}\bigr).
\end{equation}
Consequently,
\begin{equation}\label{eq:fiber-second-moment}
\sum_{x\in X_m} d_m(x)^2
=
\frac{1}{F_{m+2}}
\sum_{j=0}^{F_{m+2}-1}
\Bigl|
\prod_{k=1}^m \bigl(1+\zeta^{jF_{k+1}}\bigr)
\Bigr|^2.
\end{equation}
\end{theorem}

\begin{proof}
By Theorem~\ref{thm:fiber-residue},
\[
d_m(x)
=
\sum_{\omega\in\{0,1\}^m}
\ind_{\{\Val_m(\omega)\equiv r_x \pmod{F_{m+2}}\}}.
\]
Use the finite Fourier inversion formula on $\mathbb{Z}/F_{m+2}\mathbb{Z}$:
\[
\ind_{\{n\equiv r_x\}}
=
\frac{1}{F_{m+2}}
\sum_{j=0}^{F_{m+2}-1} \zeta^{j(n-r_x)}.
\]
Therefore
\[
d_m(x)
=
\frac{1}{F_{m+2}}
\sum_{j=0}^{F_{m+2}-1}
\zeta^{-jr_x}
\sum_{\omega\in\{0,1\}^m}
\zeta^{j\Val_m(\omega)}.
\]
Since
\[
\Val_m(\omega)=\sum_{k=1}^m \omega_k F_{k+1},
\]
the inner sum factorizes:
\[
\sum_{\omega\in\{0,1\}^m}\zeta^{j\Val_m(\omega)}
=
\prod_{k=1}^m \sum_{\omega_k\in\{0,1\}}\zeta^{j\omega_k F_{k+1}}
=
\prod_{k=1}^m (1+\zeta^{jF_{k+1}}).
\]
This proves \eqref{eq:fiber-dft}.

For \eqref{eq:fiber-second-moment}, square \eqref{eq:fiber-dft}, sum over $x\in X_m$, and use that $r_x$ runs through the complete residue system $\{0,\dots,F_{m+2}-1\}$ by Proposition~\ref{prop:finite-window-range}. Orthogonality of characters collapses the double sum.
\end{proof}

\begin{remark}
Theorem~\ref{thm:fiber-dft} makes $d_m(x)$ an explicit arithmetic object: fiber cardinalities enter the entropy decomposition of Section~\ref{sec:entropy-ledger} and the uniform collision statistics of Proposition~\ref{prop:uniform-collision} quantitatively. Figure~\ref{fig:fiber-distribution} displays the resulting fiber-size distributions for several values of~$m$.
\end{remark}

\begin{figure}[t]
\centering
\includegraphics[width=\textwidth]{figures_jpa/fig_fiber_distribution.pdf}
\caption{Fiber-size distributions for the stabilization map $\Fold_m$ at window lengths $m=5,8,12$, obtained by exact enumeration over all $2^m$ binary words. Each bar gives the number of residue classes in $X_m$ whose fiber has a given size $d_m(x)$. The distributions are highly non-uniform. They encode the combinatorial content of the coarse-graining: both the entropy ledger of Theorem~\ref{thm:kl-ledger} and the collision probability of Proposition~\ref{prop:uniform-collision} are determined by this fiber spectrum.}
\label{fig:fiber-distribution}
\end{figure}

\subsection{Affine transfer to a shifted-weight residue model}\label{sec:stab:transfer}

For auxiliary arithmetic arguments it is convenient to compare the canonical valuation
\[
\Val_m(\omega)=\sum_{j=1}^m \omega_j F_{j+1}
\]
with the shifted-weight residue model
\[
\widetilde{\Val}_m(\omega):=\sum_{j=1}^m \omega_j F_j
\qquad\text{in }\ZZ/F_{m+2}\ZZ.
\]
By Proposition~\ref{prop:finite-window-range}, every residue $r\in\{0,\dots,F_{m+2}-1\}$ has a unique canonical representative
\[
x_r:=\bigl(\Val_m|_{X_m}\bigr)^{-1}(r)\in X_m.
\]
Accordingly we write
\[
d_m^\sharp(r):=d_m(x_r)
\]
for the canonical fiber size indexed by residue, and
\[
\widetilde d_m(r):=\#\bigl\{\omega\in\{0,1\}^m:\widetilde{\Val}_m(\omega)\equiv r \pmod{F_{m+2}}\bigr\}
\]
for the shifted-weight fiber count.

\begin{lemma}[Fibonacci cross-multiplication modulo $F_{m+2}$]\label{lem:cross-multiplication}
Extend the Fibonacci sequence by $F_0:=0$. Then for every $m\ge 1$ and $1\le j\le m$,
\[
F_{m+1}F_j \equiv (-1)^{j-1}F_{m+2-j}\pmod{F_{m+2}}.
\]
\end{lemma}

\begin{proof}
Use d'Ocagne's identity
\[
F_uF_{v+1}-F_{u+1}F_v = (-1)^vF_{u-v}
\]
with $(u,v)=(m+1,j-1)$. This gives
\[
F_{m+1}F_j-F_{m+2}F_{j-1}=(-1)^{j-1}F_{m+2-j}.
\]
Reducing modulo $F_{m+2}$ proves the claim.
\end{proof}

\begin{theorem}[Affine transfer principle]\label{thm:affine-transfer}
For every $m\ge 1$, define
\[
C_m:=\sum_{\substack{1\le k\le m\\k\not\equiv m\!\!\!\pmod 2}} F_{k+1}
\]
and the affine residue map
\[
\pi_m(r):=F_{m+1}r+C_m \pmod{F_{m+2}}.
\]
Then $\pi_m$ is a bijection of $\ZZ/F_{m+2}\ZZ$, and
\[
\widetilde d_m(r)=d_m^\sharp\bigl(\pi_m(r)\bigr)\qquad\text{for all }r\in\ZZ/F_{m+2}\ZZ.
\]
In particular, the canonical and shifted-weight models have identical fiber multisets.
\end{theorem}

\begin{proof}
Define a map $\tau_m:\{0,1\}^m\to\{0,1\}^m$ by
\[
\tau_m(\omega)_k :=
\begin{cases}
\omega_{m+1-k}, & k\equiv m\!\!\!\pmod 2,\\[1mm]
1-\omega_{m+1-k}, & k\not\equiv m\!\!\!\pmod 2.
\end{cases}
\]
It is a composition of bit reversal with parity-selective bit flips, hence a bijection.

We now compute $\Val_m(\tau_m(\omega))$ modulo $F_{m+2}$. By the definition of $\tau_m$,
\[
\Val_m(\tau_m(\omega))
=
\sum_{\substack{1\le k\le m\\k\equiv m\!\!\!\pmod 2}} \omega_{m+1-k}F_{k+1}
{}+\sum_{\substack{1\le k\le m\\k\not\equiv m\!\!\!\pmod 2}} (1-\omega_{m+1-k})F_{k+1}.
\]
Separating the constant part gives
\[
\Val_m(\tau_m(\omega))
=
\sum_{k=1}^m (-1)^{m-k}\omega_{m+1-k}F_{k+1}+C_m.
\]
Re-index with $j=m+1-k$:
\[
\Val_m(\tau_m(\omega))
=
\sum_{j=1}^m (-1)^{j-1}\omega_jF_{m+2-j}+C_m.
\]
Applying Lemma~\ref{lem:cross-multiplication} termwise to $\widetilde{\Val}_m(\omega)=\sum_{j=1}^m \omega_jF_j$ yields
\[
F_{m+1}\widetilde{\Val}_m(\omega)
\equiv
\sum_{j=1}^m (-1)^{j-1}\omega_jF_{m+2-j}
\pmod{F_{m+2}}.
\]
Comparing with the previous display gives the intertwining identity
\[
\Val_m(\tau_m(\omega))\equiv F_{m+1}\widetilde{\Val}_m(\omega)+C_m
=
\pi_m\bigl(\widetilde{\Val}_m(\omega)\bigr)
\pmod{F_{m+2}}.
\]

Since consecutive Fibonacci numbers are coprime, $\gcd(F_{m+1},F_{m+2})=1$, so multiplication by $F_{m+1}$ is invertible modulo $F_{m+2}$ and $\pi_m$ is a bijection. Therefore, for every residue $r$,
\begin{align*}
\widetilde d_m(r)
&=
\#\bigl\{\omega:\widetilde{\Val}_m(\omega)\equiv r \pmod{F_{m+2}}\bigr\}\\
&=
\#\bigl\{\omega:\Val_m(\tau_m(\omega))\equiv \pi_m(r) \pmod{F_{m+2}}\bigr\}\\
&=
\#\bigl\{u:\Val_m(u)\equiv \pi_m(r) \pmod{F_{m+2}}\bigr\}\\
&=
d_m^\sharp\bigl(\pi_m(r)\bigr),
\end{align*}
where the third equality uses that $\tau_m$ is a bijection and the last equality is Theorem~\ref{thm:fiber-residue} written in residue coordinates. This proves the theorem.
\end{proof}

\begin{corollary}[Transfer of multiset-based statistics]\label{cor:affine-transfer-stats}
For every $m\ge 1$, the shifted-weight model and the canonical model have the same multiset of fiber sizes. In particular they have the same:
\begin{enumerate}[label=(\roman*),leftmargin=*,itemsep=2pt]
  \item maximal fiber size,
  \item maximal-fiber multiplicity,
  \item second moment $\sum_x d_m(x)^2$,
  \item every statistic that depends only on the unordered fiber multiset.
\end{enumerate}
\end{corollary}

\begin{proof}
Theorem~\ref{thm:affine-transfer} identifies the two residue-count functions up to the permutation $\pi_m$.
\end{proof}

\begin{remark}[Closed form of the affine offset]
The parity sum in Theorem~\ref{thm:affine-transfer} simplifies to
\[
C_m=F_{m+1}-1.
\]
Indeed, for $m=2q$ one has
\[
C_{2q}=F_2+F_4+\cdots+F_{2q}=F_{2q+1}-1=F_{m+1}-1,
\]
and for $m=2q+1$ one has
\[
C_{2q+1}=F_3+F_5+\cdots+F_{2q+1}=F_{2q+2}-1=F_{m+1}-1.
\]
Thus the affine transfer may also be written as
\[
\pi_m(r)=F_{m+1}(r+1)-1 \pmod{F_{m+2}}.
\]
\end{remark}

\begin{proposition}[Affine reflection symmetry of the canonical fiber spectrum]\label{prop:fiber-reflection}
Let
\[
B_m:=\sum_{k=1}^m F_{k+1}=F_{m+3}-2.
\]
Then for every residue $r\in \ZZ/F_{m+2}\ZZ$,
\[
d_m^\sharp(r)=d_m^\sharp(B_m-r).
\]
Equivalently, the canonical residue-indexed fiber spectrum is invariant under the affine reflection
\[
r\longmapsto B_m-r
\qquad\text{in }\ZZ/F_{m+2}\ZZ.
\]
\end{proposition}

\begin{proof}
Consider the bitwise complement involution
\[
\iota:\{0,1\}^m\to\{0,1\}^m,\qquad
\iota(\omega)_k:=1-\omega_k.
\]
It is a bijection, and
\[
\Val_m(\iota(\omega))
=
\sum_{k=1}^m (1-\omega_k)F_{k+1}
=
\sum_{k=1}^m F_{k+1}-\Val_m(\omega)
\equiv
B_m-\Val_m(\omega)
\pmod{F_{m+2}}.
\]
Therefore $\iota$ sends the preimage of residue $r$ bijectively onto the preimage of residue $B_m-r$. By the definition of $d_m^\sharp$, this proves
\[
d_m^\sharp(r)=d_m^\sharp(B_m-r).
\]
\end{proof}

\begin{corollary}[Symmetry of the uniform visible law]\label{cor:uniform-visible-symmetry}
Let $\mu_m^{\mathrm{unif}}$ be the uniform law on $\{0,1\}^m$ and let
\[
\pi_m^{\mathrm{unif}}:=(\Fold_m)_*\mu_m^{\mathrm{unif}}
\]
be the induced visible law on $X_m$. Then
\[
\pi_m^{\mathrm{unif}}(x_r)=\pi_m^{\mathrm{unif}}(x_{B_m-r})
\qquad\text{for every }r\in\ZZ/F_{m+2}\ZZ.
\]
In particular, every odd centered residue moment of the uniform visible law vanishes.
\end{corollary}

\begin{proof}
Directly from the definition of the pushforward,
\[
\pi_m^{\mathrm{unif}}(x_r)=\frac{d_m^\sharp(r)}{2^m}.
\]
The claim now follows from Proposition~\ref{prop:fiber-reflection}.
\end{proof}

A related truncation-versus-stabilization effect can also be studied at finite range, but it is not needed here.

